%% kapa: exam=1 round=28 year=2017 period=1 form=A subject=civil_law subject_ko=민법 seq=37 number=37
%% source: Q-Net 감정평가사 기출문제 (cst00309) artlSeq=5199603 file=extracted/1008/5199603/2017년 제28회 감정평가사 1차시험 문제지/1교시 A형.pdf
%% conversion: kordoc 4.24.51 -> tools/gen1.py (auto); review=proofread-sonnet; flags=none
\begin{question}{37}
질권에 관한 설명으로 옳지 않은 것은?
\begin{choices}
  \item 채권을 질권의 목적으로 하는 경우에 채권증서가 있는 때에는 질권의 설정은 그 증서를 질권자에게 교부함으로써 그 효력이 생긴다.
  \item 질권의 목적인 채권의 변제기가 질권자의 채권의 변제기보다 먼저 도래한 경우 질권자는 제3채무자에 대하여 자신에게 변제할 것을 청구할 수 있다.
  \item 저당권으로 담보한 채권을 질권의 목적으로 한 때에는 그 저당권등기에 질권의 부기등기를 하여야 그 효력이 저당권에 미친다.
  \item 질권자는 질권의 실행방법으로서 질권의 목적이 된 채권을 직접 청구할 수 있다.
  \item 양도할 수 없는 동산은 질권의 목적이 될 수 없다.
\end{choices}
\end{question}
\answer{2}
